\chapter{Fazit und Ausblick}\label{ch:fazit}

\section{Fazit}
Das Projekt hat aus einem Scrapy-Prototyp eine skalierbare Erhebungs- und Klassifikationsplattform auf AWS gemacht und sie an einem realen, großen Anwendungsfall erprobt: \zDocsTotal{} PDFs von \zUnisTotal{} Hochschulen, zwei Großläufe mit zusammen rund 18~Stunden Laufzeit, \zSpeicherGB{}\,GB Daten, eine Endliste von \zFinalListed{} Dokumenten mit geschätzt \zEstTrue{} echten Modulhandbüchern und einer Abdeckung von \zUnisLikelyPct\,\% der Hochschulen. Der Aufbau ist modular genug, dass ein anderer Dokumenttyp ein neues Profil und ein neues Modell braucht, nicht aber einen neuen Crawler. Das ist die Brücke zu KIWIT.

\Cref{tab:akzeptanz} bewertet die Akzeptanzkriterien aus Kapitel~\ref{ch:einleitung} gegen das Ergebnis.

\begin{table}[htbp]
\centering\small
\caption{Akzeptanzkriterien und Ergebnis.}\label{tab:akzeptanz}
\begin{tabularx}{\linewidth}{@{}p{2.8cm} p{3.0cm} X@{}}
\toprule
\textbf{Kriterium} & \textbf{Erfüllt?} & \textbf{Begründung} \\
\midrule
A1 vollständig durchsuchen & weitgehend & Alle \zUnisTotal{} Hochschulen wurden bearbeitet; Seiten mit JavaScript, Login oder anderer Domain nicht erreichbar. \\
A2 alle Handbücher finden (Recall) & \textbf{teilweise} & Etwa 40--47\,\% des Erwartungswerts; 39 Hochschulen ohne Fund, 117 mit Dokumenten, aber ohne bestätigtes Handbuch. \\
A3 Nicht-Handbücher erkennen (Precision) & ja, durch LLM-Stufe & TF-IDF allein reicht nicht (25--82\,\%); die streng bestätigte Liste ist vom LLM geprüft, der Rest geschätzt $\approx$ 89\,\%. \\
A4 strukturiert speichern & ja & S3-Ablage mit Manifest (Score, Quelle, Lauf, LLM-Begründung); Transparenz je Entscheidung. Eine Datenbank für Inhalte ist offen. \\
\bottomrule
\end{tabularx}
\end{table}

Das wichtigste inhaltliche Ergebnis ist \emph{negativ} und nützlich: Der Klassifikator, so gut er auf dem Papier war, genügte als alleiniger Filter nicht. Die Kombination aus billigem Massenfilter und LLM-Zweitstufe war die Antwort, und die Ehrlichkeit der Evaluation (gruppiert nach Hochschule, Pilotstichprobe, korrigierter Auswertungsfehler) hat dieses Ergebnis überhaupt erst sichtbar gemacht. Das wichtigste betriebliche Ergebnis: Eine befristete Sandbox verlangt, dass Aufbau und Daten wiederherstellbar sind.

\section{Ausblick}
\subsection{Qualität und Recall}
\begin{itemize}
  \item \textbf{Schwelle und Modell.} Im Live-Crawl die vorgeschlagene Schwelle (0,8) statt 0,5 verwenden; den Dateinamen als Merkmal aufnehmen (Abschnitt~\ref{ch:klassifikation}); mit den rund 27.700 LLM-Labels als \emph{Hard Negatives} nachtrainieren. Die dazu nötige Retrain-Strecke aus Manifesten (\texttt{ingest --manifest} mit S3-Abruf) ist vorhanden.
  \item \textbf{Manuelle Referenz.} Die vorbereitete Stichprobe von 120 Dokumenten auswerten, um die LLM-Genauigkeit zu messen.
  \item \textbf{Verdächtige Hosts.} Die ungeprüften Positive der 159 Hosts mit niedriger Bestätigungsrate LLM-prüfen (rund 15--20\,\$).
  \item \textbf{Seitenebene statt Dateiebene.} JavaScript-Portale (z.\,B.\ IU) liefern kein PDF; die Klassifikation von gerenderten Katalogseiten ist gebaut und wartet auf einen Einsatz mit Playwright oder Firecrawl.
  \item \textbf{Manuelle Nachsuche} für die 39 leeren und die kleinen privaten Hochschulen, deren Handbücher nicht online liegen oder hinter Login.
\end{itemize}

\subsection{Inhaltliche Erschließung}
Aus „ist ein Modulhandbuch“ soll „was steht darin“ werden: Module, ECTS-Punkte, Lernziele und Prüfungsformen als strukturierte, vergleichbare Daten. Für diese Extraktion ist ein LLM fachlich naheliegend; erst damit lassen sich Studiengänge zählen und systematisch vergleichen, wie es der ursprüngliche Anwendungsfall der FH Wedel verlangt.

\subsection{Plattform und Betrieb}
\begin{itemize}
  \item Modell-Artefakte in S3 mit Registry versionieren, sodass Modellwechsel ohne Neubau des Images möglich ist.
  \item Scoring vom Crawl entkoppeln (S3-Ereignis lösen Klassifikations-Lambda aus) und Jobstatus in einer Tabelle halten (Roadmap Phase~3--4).
  \item Secrets Manager statt \texttt{.env}-Dateien, echtes TLS-Zertifikat statt selbstsigniertem, SQLite $\to$ RDS bei mehreren Schreibern, Auto-Scaling hinter einem Load Balancer.
  \item Automatisierte Tests und Monitoring (Fehlerraten, Review-Band-Anteil als Drift-Indikator).
  \item Webapp nur bei Bedarf betreiben: Die dauerhaft laufende Instanz kostete im Budget mehr als jeder Crawl; ein On-Demand-Betrieb wäre günstiger.
\end{itemize}

\subsection{Anschluss an KIWIT}
Für Forschungsprofil~A der KIWIT-Gruppe (Makroebene) lassen sich damit künftig auch KI-bezogene Strategien, Richtlinien und Angebote erheben: Dazu genügen ein neues Extraktionsprofil mit Suchbegriffen, eine neue Spezifikation für die Zweitstufe und ein passendes Trainingsset. Die in diesem Projekt gewonnene Erfahrung, dass \emph{Auffinden} und \emph{Aussortieren} getrennte Probleme sind und beide eine eigene Lösung brauchen, sollte dabei von Beginn an berücksichtigt werden.
